% Clause 6: source PDF page 34 / printed page 22.
% Fidelity pass 03: each rule statement is followed by an additional 3.3 pt
% source gap.  This accumulation reproduces the page's vertical rule cadence.
% The page is enlarged locally so the source's final 707.8 pt baseline fits
% without moving the calibrated running footer.

\newpage
\enlargethispage{18pt}
\clause{6. Federate rules}

This clause describes the five rules that apply to HLA federates. Rationale for the inclusion of each rule is\linebreak
located in Annex A.

\subclause{6.1 Rule 6}

\textit{Federates shall have an HLA SOM, documented in accordance with the HLA OMT.}
\vspace{3.3pt}

The HLA SOM shall include the object classes, class attributes, and interaction classes of the federate that\linebreak
can be made public in a federation. The SOM is specified using one or more SOM modules and one optional\linebreak
MIM. The HLA does not prescribe which data are included in the SOM; this task shall be the responsibility\linebreak
of the federate developer. SOMs shall be documented in the format prescribed in IEEE Std 1516.2-2010.

\subclause{6.2 Rule 7}

\textit{Federates shall be able to update and/or reflect any instance attributes and send and/or receive interactions,\linebreak
as specified in their SOMs.}
\vspace{3.3pt}

The HLA allows for joined federates to make internal object representations and interactions available for\linebreak
external use as part of federation executions. These capabilities for external interaction shall be documented\linebreak
in the SOM for the federate. If documented in the SOM, these federate capabilities shall include the\linebreak
obligation to export updated values of instance attributes that are calculated internally in the federate and the\linebreak
obligation to be able to exercise interactions represented externally (i.e., by other federates in a federation).

\subclause{6.3 Rule 8}

\textit{Federates shall be able to transfer and/or accept ownership of instance attributes dynamically during a\linebreak
federation execution, as specified in their SOMs.}
\vspace{3.3pt}

The HLA allows ownership of instance attributes of an object instance to be transferred dynamically during\linebreak
a federation execution. The instance attributes of a federate that can be either owned or reflected, and whose\linebreak
ownership can be dynamically acquired or divested during execution, shall be documented in the SOM for\linebreak
that federate.

\subclause{6.4 Rule 9}

\textit{Federates shall be able to vary the conditions (e.g., thresholds) under which they provide updates of\linebreak
instance attributes, as specified in their SOMs.}
\vspace{3.3pt}

The HLA permits federates to own (i.e., provides the privilege to produce updated values for) instance\linebreak
attributes of object instances represented in the federate and to then make those values available to other\linebreak
federates through the RTI. Different federations may specify different conditions under which instance\linebreak
attributes will be updated [e.g., at some specified rate, or when the amount of change in value exceeds a\linebreak
specified threshold (such as altitude changes of more than 1000 ft)]. The conditions applicable to the update\linebreak
of specific instance attributes owned by a federate shall be documented in the SOM for that federate.

\subclause{6.5 Rule 10}

\textit{Federates shall be able to manage local time in a way that will allow them to coordinate data exchange with\linebreak
other members of a federation.}
\vspace{3.3pt}

Federation designers will identify their time management approach as part of their implementation design.\linebreak
Federates shall adhere to the time management approach of the federation.
